% R1 — within-COVID mortality discrimination (synced to P0 round-17 state, 2026-09-09)
\subsection{Within-COVID, models differ from each other, but no foundation model beats the from-scratch baseline under the pre-specified test}
\label{ssec:r1_within_covid}

We first asked whether foundation-model embeddings carry any patient-level
mortality signal within COVID-19. For each model, frozen embeddings of the same
raw platelet transcriptomes were classified with logistic regression under
repeated patient-level cross-validation (10 repeats of 5-fold
\texttt{StratifiedGroupKFold} with \texttt{patient\_id} as the group), and each
patient's out-of-fold cell predictions were averaged into a single score before
computing the patient-level out-of-fold ROC-AUC (patient-level AUC;
Section~\ref{ssec:methods_metrics}). Three from-scratch controls ran under the
identical protocol: a PCA-200 + logistic-regression baseline on raw counts, a
null model given only each patient's cell count and disease label, and a
cohort-only model given each patient's cohort mortality rate.

Patient-level mortality signal was learnable: five of the six methods,
including the from-scratch baseline, significantly outscored the null model
(exact DeLong paired tests on the seed-ensemble scores, $p = 0.0006$--$0.023$);
the exception was \model{UCE} ($p = 0.157$; Supplementary
Note~\ref{sn:method_vs_null}). The pre-specified primary contrast
(Section~\ref{ssec:methods_cv}) asked whether the best foundation model beats
the from-scratch baseline once winner selection is accounted for. It was not
met: the observed maximum advantage (\model{Geneformer}, $+0.064$
patient-level AUC on the seed-ensemble scores) did not exceed what the
selection procedure produces under the null (patient-level label-permutation
max-statistic test, 10{,}000 permutations, $P = 0.181$; recentered-bootstrap
cross-check, $P = 0.073$; Supplementary Note~\ref{sn:winner_test}). The result
is stable: all 25 leave-one-fatal-patient-out refits returned $P > 0.05$
(range 0.141--0.288; Supplementary Note~\ref{sn:fragility}).

The model ranking underneath this null is still informative
(Table~\ref{tab:within_covid} and Figure~\ref{fig:within_covid}). \model{Geneformer} and \model{STATE} led at
0.786 $\pm$ 0.026 and 0.778 $\pm$ 0.033 (repeat-mean $\pm$ repeat-sd) and
could not be distinguished from each other ($\Delta = -0.012$, exact DeLong
$p = 0.58$). Their nominal advantages over the baseline (\model{Geneformer}
$+0.064$, $p = 0.036$; \model{STATE} $+0.052$, $p = 0.048$) are reported as
exploratory and do not survive the winner correction above. All remaining
contrasts are likewise exploratory
(Supplementary Note~\ref{sn:se_check}):
\model{scPRINT-2} placed third (0.740 $\pm$ 0.027) and could not be separated
from the baseline ($p = 0.65$), and \model{scGPT} (0.720 $\pm$ 0.047) was
statistically indistinguishable from it ($p = 0.91$). \model{UCE} placed last
(0.661 $\pm$ 0.046), below the from-scratch baseline, though the paired
contrast was not itself significant ($-0.066$, $p = 0.081$). One separation
survives selection correction: the two leaders' advantage over \model{UCE}
($+0.118$ and $+0.130$) is confirmed by a range-statistic test on the same
label permutations as the primary analysis (max-minus-min over the five
foundation models, $P = 0.011$; Supplementary Note~\ref{sn:range_test}) and
is the paper's single confirmatory secondary result. Within the exploratory
pairwise family, six contrasts survive Benjamini--Hochberg correction, all
with \model{UCE} or \model{scGPT} on the losing side ($q \leq 0.037$;
Supplementary Figure~\ref{fig:supp_pairwise}), and the two
leaders-versus-\model{UCE} gaps are the largest. The result's content is a
statement about the losing model, addressed in
Section~\ref{ssec:discussion_property}.

A separate question, whether a model's score adds information beyond the
baseline's even without ranking patients better on its own, was examined with
nested-model comparisons: a reduced model containing cohort and the PCA-200
score against a full model adding each foundation-model score in turn. No
model's increment survived the same Benjamini--Hochberg correction used
throughout, in any of three comparator configurations
(Supplementary Note~\ref{sn:nested_comparison}). The complementarity question
therefore returns the same answer as the superiority question, and the
headline statement needs no qualifier.

Three controls bound where the learnable signal does not come from. The null
model reached only 0.554, so mortality was not readable from cell counts or
disease identity alone; per-patient cell counts also correlated only weakly
with predicted scores across models (Spearman $\rho = 0.18$--$0.29$;
Supplementary Figure~\ref{fig:supp_cellcount}). The cohort-only model reached
0.486 $\pm$ 0.038, statistically indistinguishable from chance; because a
rate estimator based on add-one smoothing inverts risk ordering for small
cohorts, this number is silent on the confound question rather than
exonerating cohort (Supplementary Note~\ref{sn:cohort_baseline}). A
QC-metrics-only model scored below chance out-of-fold (0.376 as-fit), a
textbook suppression pattern: the QC features carry outcome information only
conditional on library depth, and adding each patient's cell count recovers a
correctly oriented 0.618 (Supplementary Note~\ref{sn:qc_baseline}). The
decisive test is incremental value over all of these covariates at once:
a baseline of QC metrics, cell count, and cohort identity reached 0.627, and
the pre-specified primary for this question, whether the from-scratch
platelet representation adds information over that baseline, chosen because
it involves no model selection, was met by a nested likelihood-ratio test
against an empirical permutation null ($P < 0.0001$;
Supplementary Note~\ref{sn:qc_incremental}). The five foundation models,
reported as exploratory with both readings paired, each added significant
value by the same test ($P < 0.0001$--$0.0022$, \model{UCE} weakest) and, for
all but \model{UCE}, by paired-AUC comparison ($\Delta = +0.026$ to $+0.065$,
exact DeLong $p = 0.0001$--$0.006$; \model{UCE} $+0.016$, $p = 0.072$). The
likelihood-ratio family establishes that incremental information exists; the
AUC increments quantify it; and the resolution argument applies to that
magnitude, not to the likelihood-ratio $P$-values.

Two protocol observations qualify the headline numbers. First, random-forest
classifiers scored uniformly lower than logistic regression on the same
embeddings (e.g., \model{STATE}: 0.696 versus 0.778), so we report logistic
regression throughout and retain random forest only for calibration analysis
(Section~\ref{ssec:r5_calibration}). Second, per-seed variation was large enough
to change the apparent ranking: \model{STATE}'s patient-level AUC ranged from
0.711 to 0.815 across the ten repeats, and a single 5-fold split would have
reported any point in that range as ``the'' result. We therefore treat the
repeat-mean, not any single split, as the primary estimate throughout, with
inferential quantities computed on the seed-ensemble scores
(Section~\ref{ssec:methods_cv}).

The within-COVID picture is thus one of inter-model differences without a
detectable pretraining advantage. Signal is present and learnable by nearly
every method, and it is not reducible to library quality, cell counts, or
cohort membership: platelet representations carry outcome information beyond
all three. Yet no foundation-model representation carries information beyond
what a from-scratch PCA of the same data already provides, not in
discrimination and not in nested complementarity. The two leading models
separate from \model{UCE} under a confirmatory test, a statement about the
weakest model rather than the strongest. At 25 fatal events this benchmark
resolves advantages of roughly 0.10 patient-level AUC or larger between two
different representations (Section~\ref{ssec:discussion_property}): the
nominal 0.05--0.06 advantages sit below that resolution, and detecting effects
of that size would require on the order of $10^3$ patients. Whether any signal
transfers beyond COVID-19 is addressed next
(Sections~\ref{ssec:r2_within_sepsis}--\ref{ssec:r3_cross_disease}).
